\documentclass[10pt,a4paper]{amsart}
\usepackage[margin=19mm]{geometry}
\usepackage{amsmath,amssymb,mathtools}
\usepackage[scr=rsfs,cal=euler]{mathalfa}
\usepackage{xfrac}
\usepackage[unicode,hidelinks,bookmarks=false]{hyperref}
\newcommand{\ie}{i.e.\ }
\newcommand{\cf}{cf.\ }
\newcommand{\resp}{respectively\ }
\newcommand{\Subsection}{\subsection}
\providecommand{\nobreakdash}{\nobreak\hbox{-}}
\setlength{\emergencystretch}{2em}
\multlinegap0pt
\usepackage{amsthm}
\usepackage{mathtools}
\usepackage{amsmath,amssymb,amsfonts}
\usepackage{graphicx}
\usepackage{bm}
\usepackage{enumitem}
\usepackage{mathrsfs}
\usepackage{xcolor}
\providecommand{\lrabs}[1]{\!\left\lvert #1 \right\lvert}
\providecommand{\lrb}[1]{\!\left[#1\right]}
\theoremstyle{plain}
\newtheorem*{unnumberedtheorem}{Theorem}
\newtheorem*{unnumberedproposition}{Proposition}
\newtheorem*{unnumberedlemma}{Lemma}
\newtheorem{theorem}{Theorem}[section]
\newtheorem{proposition}[theorem]{Proposition}
\newtheorem{lemma}[theorem]{Lemma}
\newtheorem{corollary}[theorem]{Corollary}
\newtheorem{conjecture}[theorem]{Conjecture}
\theoremstyle{definition}
\newtheorem{remark}[theorem]{Remark}
\numberwithin{equation}{section}
\numberwithin{figure}{section}
\numberwithin{table}{section}
\title[On murmurations]{On Murmurations: the Lipschitz bound for the Chebyshev-type polynomial at a fixed base point (Lemma 2.9)}
\author[T. Nelson, E. Ross, M. Wassercug, H. Xue]{Timothy Nelson, Erick Ross, Maya Wassercug, Hui Xue}
\date{Source: arXiv:2609.29631v1 (CC BY 4.0)}
\begin{document}
\maketitle
\vspace{1em}
\noindent\textit{Source and licence.} On murmurations. Version arXiv:2609.29631v1 (CC BY 4.0), distributed under the Creative Commons Attribution 4.0 International licence, http://creativecommons.org/licenses/by/4.0/, as recorded in the arXiv OAI licence metadata for this exact version and shown on the abstract page https://arxiv.org/abs/2609.29631v1. This item is an editorial re-presentation of the paper's uniform Lipschitz bound for the Chebyshev-type polynomial of the second kind at a fixed base point, printed as Lemma 2.9. A full rights notice is delivered at licenses/arxiv-2609.29631-CC-BY-4.0.txt.
\noindent\textit{Editorial scope.} Counted unit: the paper's uniform Lipschitz bound for the Chebyshev-type polynomial of the second kind at a fixed base point, printed as Lemma 2.9 (source0.tex lines the statement at lines 709--715 and the proof at lines 716--801), delivered together with the paper's complete original proof of it as the single primary proof body. No mathematics is written, repaired or re-derived: statement and proof are the paper's own text line for line, in the paper's own notation. Every cross-reference inside the retained spans resolves inside the delivered document. Printed numbers are the paper's own: the wrapper restores the paper's counters before the retained material, so the counted statement prints as Lemma 2.9, the paper's own number for it, and its display carries the paper's own equation number. Omitted from this package and not redistributed: the paper's preamble and front matter as such (of which only the title and the author names are reproduced in the wrapper); the sections, lemmas, theorems and examples that lie outside the retained spans; and the paper's bibliography with its bib data file. No citation command occurs in any retained span, so the delivered document carries no bibliography and no reference or citation command of the delivered text was rewritten or adapted. Printed slips kept literal, among them the paper's own wording and its own choice of symbols. Layout and class: the delivered wrapper is the lane's pinned amsart editorial wrapper at 10pt on A4, to which this package adds the paper's own theorem-like declarations and the shorthand macros its retained text uses, all copied verbatim from the paper's source. No publisher class file, no publisher style file and no upstream script is redistributed. Assets: no retained span of this package contains a figure environment or a graphics-inclusion command, no image or artwork file exists in or is redistributed with this package, and the manifest and the delivery evidence both carry an empty asset-dependency list. A full rights notice is delivered at licenses/arxiv-2609.29631-CC-BY-4.0.txt.
\setcounter{section}{2}
\setcounter{theorem}{8}
\begin{lemma} \label{lem:lipschitz-const-at-x0-for-U}
    Fix a natural number $K \ge 1$, and let $x_0 \in (-1,1)$. Then for all $x \in [-1,1]$,
    \begin{align}
        \label{eqn:lipschitz-const-goal}
        |U_{K-1}(x) - U_{K-1}(x_0)| \le \frac{3K}{1-|x_0|} |x-x_0|.
    \end{align}
\end{lemma}

\begin{proof}
We divide into two cases.

\textbf{Case 1: $|x-x_0| > \frac{2}{3} (1-|x_0|)$. \\}
Recall the global bound for Chebyshev polynomials
\begin{align}
    -K \le U_{K-1}(x) \le K.
\end{align}
This bound then implies that
\begin{align}
    |U_{K-1}(x)-U_{K-1}(x_0)| \le 2K = \frac{3K}{1-|x_0|} \cdot \frac{2}{3} (1-|x_0|) \le \frac{3K}{1-|x_0|} |x-x_0|,
\end{align}
verifying \eqref{eqn:lipschitz-const-goal}.


\textbf{Case 2: $|x-x_0| \le \frac{2}{3} (1-|x_0|)$. \\}
Observe that for every $t$ between $x_0$ and $x$, 
\begin{align}
    1-|t|
    &\ge
    1-|x_0|-|t-x_0| \\
    &\ge
    1-|x_0|-\frac{2}{3}(1-|x_0|) \\
    &=
    \frac{1}{3}(1-|x_0|).
    \label{eq:xi-away-from-endpoint}
\end{align}


Additionally, we have the derivative bound
\begin{align}
    |U_{K-1}'(t)|
    &=
    \lrabs{
        \frac{1}{-\sin \theta}
        \cdot
        \frac{d}{d\theta}
        \lrb{
            \frac{\sin(K\theta)}{\sin \theta}
        }
    } \qquad \qquad\quad
    \text{(where $t=\cos \theta$)}
    \\
    &=
    \lrabs{
    \frac{
        K\cos(K\theta)\sin\theta
        -
        \sin(K\theta)\cos\theta
    }{\sin^3\theta}
    } \\
    &\le
    \frac{
        K\sin\theta
        +
        K\sin\theta\,|\cos\theta|
    }{\sin^3\theta}
    \qquad\qquad\quad
    \text{(since $|\sin(K\theta)|\le K\sin\theta$)}
    \\
    &=
    \frac{K(1+|t|)}{1-t^2} \\
    &=
    \frac{K}{1-|t|}.
    \label{eq:U-derivative-bound}
\end{align}



Hence by the Mean Value Theorem, we have that
\begin{align}
    |U_{K-1}(x)-U_{K-1}(x_0)|
    &=
    |U_{K-1}'(t)|\,|x-x_0| 
    \qquad \text{(for some $t$ between $x_0$ and $x$)} 
    \\
    &\le
    \frac{K}{1-|t|}|x-x_0| 
    \qquad\quad\ \text{(by \eqref{eq:U-derivative-bound})}
    \\
    &\le
    \frac{3K}{1-|x_0|}|x-x_0|, 
    \qquad\ \,  \text{(by \eqref{eq:xi-away-from-endpoint})}
\end{align}
verifying \eqref{eqn:lipschitz-const-goal}.
\end{proof}

\end{document}
